Investigar seis aplicaciones nativas del sistema operativo y su importancia en la gestión de seguridad informática.

Kali Linux es una distribución de Linux basada en Debian que tiene algunas herramientas instaladas por defecto que se enfocan en las tareas de seguridad como lo son el \textit{Penetration Testing}, investigación de seguridad, computación forense e ingeniería inversa~\cite{kali_docs}.

Algunas de las herramientas nativas del sistema con las que podemos realizar tareas de seguridad son:
\begin{itemizeColor}
    \item \textbf{Aircrack-ng:} Es una suite completa de herramientas que nos permite trabajar sobre la seguridad inalámbrica. La herramienta nos permite monitorear paquetes para su posterior procesamiento con herramientas de terceros, realizar inyección de paquetes para realizar un ataque, probar las capacidades de los controladores y las tarjetas Wi-Fi, así como descifrar contraseñas WEP y WPA/WPA2-PSK. Al ser la red inalámbrica uno de los medios por los que la mayoría de usuarios se comunican, esta herramienta puede ser muy valiosa para encontrar vulnerabilidades dentro de la red que puedan causar ataques, sobre todo dentro de redes públicas~\cite{aircrack_ng_docs}.
    \item \textbf{Hydra:} Es una herramienta de craqueo de contraseñas de inicios de sesión de red que permite numerosos protocolos de ataque, lo que permite a los consultores de seguridad el demostrar qué tan fácil puede ser obtener un acceso no autorizado a un sistema de forma remota. La herramienta de craqueo de contraseñas nos permite saber qué tan robustas son las credenciales de acceso~\cite{hydra_github}.
    \item \textbf{John the Ripper:} Es un descifrador de contraseñas rápido que nos permite recuperar contraseñas para poder detectar contraseñas débiles mediante ataques con hashes criptográficos. Esta herramienta nos permite prevenir que las contraseñas sean explotadas por atacantes, así como comprobar la solidez de las credenciales en el pentesting~\cite{john}.
    \item \textbf{Nmap:} Es una herramienta para la exploración de red y auditoría de seguridad que nos permite, entre otras cosas, determinar qué equipos se encuentran disponibles en una red, qué servicios ofrecen, qué sistemas operativos ejecutan, los filtros de paquetes o cortafuegos que utilizan, así como docenas de otras características. Esta herramienta nos permite hacer auditorías de seguridad de la red al identificar información relevante de los dispositivos conectados~\cite{nmap}.
    \item \textbf{Sqlmap:} Es una herramienta que automatiza la detección y explotación de vulnerabilidades de inyección SQL en aplicaciones web para la toma de control de servidores de bases de datos. Nos permite realizar pentesting para ayudar a identificar fallas en sistemas web que puedan ser explotadas por atacantes~\cite{sqlmap}.
    \item \textbf{Wireshark:} Esta herramienta es un sniffer o analizador de protocolos de red que permite identificar anomalías al permitirnos inspeccionar en tiempo real el tráfico de una red. Es útil para identificar actividades sospechosas en red, así como ver qué ocurrió después de un ataque, por lo que se utiliza en auditorías de seguridad de redes para verificar que no haya fugas de información~\cite{wireshark}.
\end{itemizeColor}
